
\subsubsection{3.1 Overview}

DNSI quantifies benchmark saturation by measuring the entropy of performance improvements over time, normalized by task difficulty:

$$\text{DNSI} = \frac{H(\Delta_{\text{observed}})}{H_{\text{expected}}(D)}$$

where $H(\Delta_{\text{observed}})$ is the entropy of observed improvement deltas and $H_{\text{expected}}(D)$ is the expected entropy given difficulty proxy $D$.

A healthy benchmark exhibits diverse improvements producing high improvement entropy. A saturated benchmark shows clustered, incremental improvements producing low entropy. Normalizing by difficulty separates true saturation from benchmark hardness.

\subsubsection{3.2 Computation Steps}

\textbf{Step 1: Extract SOTA History.} From PapersWithCode or equivalent repositories, extract the time series of SOTA performance:

$$S = \{(t_1, p_1), (t_2, p_2), \ldots, (t_n, p_n)\}$$

where $t_i$ is the submission timestamp and $p_i$ is the reported performance.

\textbf{Step 2: Compute Improvement Deltas.} Compute windowed improvement deltas using 6-month aggregation:

$$\Delta_w = \sum_{t_i \in w} (p_i - p_{i-1})^+$$

where $w$ indexes 6-month windows and $(x)^+ = \max(0, x)$ counts only improvements. The 6-month window smooths conference clustering effects.

\textbf{Step 3: Compute Improvement Entropy.} Discretize deltas into bins and compute Shannon entropy:

$$H(\Delta) = -\sum_{b} p_b \log_2 p_b$$

where $p_b$ is the proportion of deltas falling in bin $b$.

\textbf{Step 4: Difficulty Normalization.} Normalize by expected entropy given task difficulty:

$$H_{\text{expected}}(D) = \log_2(D)$$

For image classification, the difficulty proxy is the number of classes: $D_{\text{vision}} = N_{\text{classes}}$.

\textbf{Final Formula:}

$$\text{DNSI} = \frac{H(\Delta_{\text{observed}})}{\log_2(N_{\text{classes}})}$$

\textbf{Interpretation:} DNSI $\approx$ 1.0 indicates improvement entropy matches expected diversity (healthy benchmark). DNSI < 0.5 indicates entropy significantly below expected (saturated benchmark). DNSI > 1.0 indicates more diversity than expected.

\subsubsection{3.3 Design Decisions}

Major ML conferences cluster submissions, creating artificial periodicity in raw SOTA histories. Six-month windows smooth this clustering while preserving the saturation signal. For classification tasks, the number of classes provides an information-theoretic bound on output complexity. The minimum requirement is >15 SOTA entries over >3 years for reliable entropy estimation.
